\documentclass[11pt]{article}
\usepackage[margin=1in]{geometry}
\usepackage{parskip}
\usepackage[hidelinks]{hyperref}
\pagestyle{empty}

\begin{document}

\noindent Phuc Hao Do \\
Center for AI Innovation, Research, and Application (CAIRA) \\
Da Nang Architecture University (DAU), Da Nang, Vietnam \\
\href{mailto:haodp.sut@gmail.com}{haodp.sut@gmail.com}

\vspace{1em}
\noindent 19 June 2026

\vspace{1em}
\noindent To the Guest Editors, \\
\textit{IEEE Network} Special Issue on ``Quantum-AI Synergies Empowering IoT Network Resilience''

\vspace{1em}
\noindent Dear Guest Editors,

I am pleased to submit the manuscript entitled \textbf{``Self-Healing IoT Networks: A Quantum-AI Closed-Loop Framework for Resilience''} for consideration in the Special Issue on Quantum-AI Synergies Empowering IoT Network Resilience.

The article speaks directly to the central theme of the special issue: how the synergy between artificial intelligence and quantum-inspired optimization can empower IoT network resilience. We argue that resilience calls for a closed healing loop rather than one-shot repair, and we present a framework in which a lightweight AI perception layer detects and localizes degradation while a quantum-inspired optimization layer recomputes a resilient configuration under an energy budget. The synergy is concrete and instrumented rather than rhetorical: the AI risk map both warm-starts and persistently prunes the optimizer's probabilistic qubit state, and recovered configurations are fed back to keep the detector calibrated.

We evaluate the loop on two structurally distinct threats, random cascading node failures and adversarial spatial link jamming, over fifty random scenarios each and against genetic-algorithm, simulated-annealing, and greedy baselines. Myopic reactive repair is shown to leave networks fragile, while the loop restores them; the AI warm-start helps significantly, with high statistical significance on the jamming instances. We report the comparison honestly: the quantum-inspired optimizer is competitive with strong metaheuristics rather than dominant, and an ablation pinpoints the operative mechanism. Importantly, we further show that the \emph{same} AI prior warm-starts a simulated quantum approximate optimization algorithm, making the bridge from our quantum-inspired method to genuine quantum hardware concrete.

The manuscript conforms to the magazine's format. The body is within the word limit, it uses six figures and tables in total, and, in keeping with the magazine's guidance, we have limited the mathematics to three simple equations, for which we respectfully request the Guest Editors' consent. It cites fifteen archival references. All code and data needed to reproduce every result are publicly available at \url{https://github.com/haodpsut/selfheal-iot-qai}.

This work is original, has not been published previously, and is not under consideration for publication elsewhere. The author has no conflicts of interest to declare.

We believe the article offers an honest, reproducible, and forward-looking contribution squarely within the scope of the special issue, and we hope you find it suitable for peer review. Thank you for your time and consideration.

\vspace{1.5em}
\noindent Sincerely, \\[1em]
Phuc Hao Do \\
CAIRA $\cdot$ Da Nang Architecture University

\end{document}
